\section*{Erd\H{o}s Problem 737: one edge in all sufficiently long cycle lengths}
\subsection*{1. Statement and proof dependency}
If $\chi(G)=\aleph_1$, must there be an edge $e$ belonging to a cycle of length $n$ for every sufficiently large integer $n$? Yes. We reconstruct Thomassen's proof and explicitly use the following older theorem of Erd\H{o}s and Hajnal:
\begin{quote}
Every graph of uncountable chromatic number contains every finite bipartite graph as a subgraph.
\end{quote}
Only its consequence that every such graph contains $K_{t,t}$ for each finite $t$ is needed. This is an external theorem, not reproved here. All choices needed to keep the same edge for every large length are proved below.
\subsection*{2. A vertex in arbitrarily large complete bipartite graphs}
In any uncountably chromatic graph $H$, some vertex $v$ lies in a copy of $K_{t,t}$ for every $t\ge1$. Otherwise assign to each vertex $u$ the least positive integer $t(u)$ for which no $K_{t(u),t(u)}$ in $H$ contains $u$. The sets $V_t=\{u:t(u)=t\}$ form a countable partition. If each $H[V_t]$ were countably colorable, the pair consisting of $t$ and its color would give a countable coloring of $H$. Therefore some $H[V_t]$ is uncountably chromatic. The external theorem produces a $K_{t,t}$ inside $H[V_t]$, contradicting the definition of every vertex in that copy.
Notice that the copies need not be nested. A single common vertex is enough.
\subsection*{3. Fixing the edge by breadth-first levels}
Some connected component of $G$ is uncountably chromatic, since otherwise countably many colors could be reused independently in all components. Work in that component and fix a root $x$. Write $L_j$ for its distance-$j$ level. Some induced graph $G[L_m]$ is uncountably chromatic: otherwise coloring by the pair (level, color within that level) gives a countable coloring. Necessarily $m\ge1$.
Choose a vertex $v\in L_m$ as in the preceding lemma, applied to $G[L_m]$. Fix a breadth-first spanning tree rooted at $x$, and let $e$ be the first edge on the tree path from $v$ toward $x$. For every distinct $y\in L_m$, the tree path $R_{v,y}$ between $v$ and $y$ has even length $2h$ with $1\le h\le m$, contains $e$, and has all its internal vertices in levels strictly below $m$. These facts follow by taking the least common ancestor of the two equal-depth vertices.
\subsection*{4. Both parities of length, with the same edge}
Fix any integer $n>2m$. Choose $t$ so large that $2t-2\ge n$, and take a $K_{t,t}$ in $G[L_m]$ containing $v$.
If $n$ is odd, choose $y$ in the opposite bipartition class from $v$. Within $K_{t,t}$ there is a simple $v$--$y$ path of every odd length $1,3,\ldots,2t-1$: alternate distinct vertices of the two classes, reserving the endpoints. Since $|R_{v,y}|=2h\le2m<n$, the integer $n-2h$ is positive, odd, and at most $2t-1$. Choose the corresponding path $Q$.
If $n$ is even, choose $y\ne v$ in the same class. The complete bipartite graph has a simple $v$--$y$ path of each even length $2,4,\ldots,2t-2$. Now $n-2h$ is a positive even integer in this range, so again choose $Q$ of this length.
In both cases $Q$ lies entirely in $L_m$, whereas the internal vertices of $R_{v,y}$ lie below $L_m$. Thus their union is a simple cycle of length $n$ containing the already fixed edge $e$. The copies, $y$, and $h$ may depend on $n$; $e$ does not.
\subsection*{5. Sources and assessment}
Sources: \url{https://www.erdosproblems.com/737}; C. Thomassen, \emph{Cycles in graphs of uncountable chromatic number}, Combinatorica 3 (1983), 133--134, \url{https://doi.org/10.1007/BF02579349}; P. Erd\H{o}s and A. Hajnal, \emph{On chromatic number of graphs and set-systems}, Acta Math. Acad. Sci. Hungar. 17 (1966), 61--99.
This is a full reconstruction relative to the explicitly stated Erd\H{o}s--Hajnal input; it is not a new resolution. The argument actually works for any uncountable chromatic number.
\textbf{SOLVED:} the single edge $e$ lies on cycles of every integer length greater than $2m$.
\textbf{Completion rate estimate: 100\%.}
